Die Anforderungen an wb-newsfeed stammen aus zwei Dokumenten.\kihinweis{}
Das erste ist das Briefing des Projektpartners Werbeberg vom 30.~Jänner 2026, das den Funktionsumfang und die technischen Rahmenbedingungen festlegt.
Das zweite ist der Diplomarbeitsantrag vom 19.~Februar 2026, in dem die geplanten Ergebnisse auf die beiden Autoren aufgeteilt sind.
Dieses Kapitel fasst beide Quellen zusammen, beschreibt die Personen, die mit dem System arbeiten, und ordnet die Anforderungen in funktionale und nicht-funktionale Anforderungen.
Jede Anforderung trägt eine Kennung, über die in späteren Kapiteln auf sie verwiesen wird.

\section{Projektauftrag und Rahmenbedingungen}
\setauthor{\firstauthor}

\reminder{Ein bis zwei Sätze zum Projektpartner Werbeberg ergänzen: Tätigkeitsfeld, Ansprechperson, warum das Thema für das Unternehmen interessant ist.}

Der Auftrag lässt sich in einem Satz zusammenfassen: Aus gewöhnlichen Nachrichtenartikeln sollen automatisch Short-Form-Inhalte entstehen, die in einem vertikal scrollbaren Feed angezeigt werden, und dieser Feed soll sich als Komponente in bestehende Nuxt-Projekte einbauen lassen.
Das Briefing beschreibt damit nicht nur eine Anwendung, sondern einen Baustein, den eine Website übernehmen kann, ohne ihr bestehendes \gls{cms} zu ersetzen.

Daraus ergeben sich drei Teilaufgaben.
Die erste ist die Generierung: Aus Titel, Text und Kategorie eines Artikels entstehen ein \gls{teaser}, Kernpunkte und weitere Metadaten, die gespeichert werden.
Die zweite ist die Bereitstellung über eine \gls{rest}-Schnittstelle, die neue und geänderte Inhalte ausliefert.
Die dritte ist die Darstellung im Feed selbst, die als wiederverwendbares Package gebaut wird.
Möchte jemand einen Artikel vollständig lesen, übergibt der Feed lediglich eine eindeutige Kennung an die einbindende Website.
Wie diese den Artikel anzeigt, bleibt ihr überlassen.

Das Briefing gibt außerdem den technischen Rahmen vor, der in Tabelle~\ref{tab:req:tech} zusammengefasst ist.
Für die Oberfläche sind Nuxt~4 mit \gls{ssr}, TypeScript und Tailwind CSS in Version~4 vorgegeben.
Für das Backend schreibt das Briefing keine Sprache vor.
Verlangt wird nur, dass es die Generierung asynchron verarbeitet und die Inhalte verwaltet und bereitstellt.

\begin{table}
    \centering
    \small
    \begin{tabularx}{\textwidth}{l >{\raggedright\arraybackslash}X}
        \toprule
        \textbf{ID} & \textbf{Rahmenbedingung} \\
        \midrule
        T1 & Oberfläche mit Nuxt~4 (\gls{ssr}) und TypeScript \\
        T2 & Gestaltung mit Tailwind CSS in Version~4 \\
        T3 & Einheitliche Code-Prüfung mit ESLint im gesamten Projekt \\
        T4 & Feed-Oberfläche als wiederverwendbares Package \\
        T5 & Backend für asynchrone Generierung, Verwaltung und Bereitstellung der Inhalte \\
        \bottomrule
    \end{tabularx}
    \caption{Technische Rahmenbedingungen laut Briefing}
    \label{tab:req:tech}
\end{table}

Die Arbeit ist entlang dieser Teilaufgaben aufgeteilt.
\firstauthor{} verantwortet das Backend mit Schnittstelle, Datenbank, Warteschlange, Text- und Bildgenerierung sowie die Dokumentation der Schnittstellen.
\secondauthor{} verantwortet die Feed-App, das Admin-\gls{cms}, das Komponenten-Package und dessen Einbindung in fremde Websites.
Die Entwicklung fand zwischen Juli und September 2026 statt.
Das gemeinsame Git-Repository umfasst in diesem Zeitraum 258 Commits.

Zur Abgrenzung wurden zu Projektbeginn einige Punkte festgelegt.
Das System ist kein Ersatz für ein Redaktionssystem, Artikel entstehen weiterhin in der Redaktion und gelangen über eine Schnittstelle oder einen Import in wb-newsfeed.
Benutzerkonten für das Publikum sind nicht vorgesehen.
Mehrere Erweiterungsideen, etwa eine Audiofassung oder Teaser-Varianten für unterschiedliche Zielgruppen, sind als optional markiert.
Sie werden in Abschnitt~\ref{sec:req:functional} getrennt von den Pflichtanforderungen geführt.
Wer mit dem System im Alltag arbeitet, beschreibt der nächste Abschnitt.

\section{Rollen und Anwendungsfälle}
\setauthor{\secondauthor}

Mit wb-newsfeed arbeiten drei Gruppen von Personen, deren Erwartungen sich deutlich unterscheiden.

\begin{compactdesc}
    \item[Leser:innen] sehen den Feed auf dem Smartphone, meist eingebettet in die Website eines Mediums. Sie haben kein Konto und erwarten, dass der Feed sofort lädt und sich wie gewohnte Short-Form-Plattformen bedienen lässt.
    \item[Redakteur:innen] arbeiten im Admin-\gls{cms}. Sie legen Artikel an oder importieren sie, prüfen die generierten Inhalte, korrigieren sie bei Bedarf und veröffentlichen die Artikel.
    \item[Entwickler:innen der Host-Website] binden den Feed in ihre bestehende Nuxt-Website ein. Für sie zählt eine überschaubare Einrichtung, eine klare Dokumentation und die Möglichkeit, das Aussehen an die eigene Marke anzupassen.
\end{compactdesc}

Die Anforderungen wurden für jede dieser Gruppen als User Stories formuliert.
Tabelle~\ref{tab:req:stories} zeigt eine Auswahl, die den Kern des Systems abdeckt.

\begin{table}
    \centering
    \small
    \begin{tabularx}{\textwidth}{l >{\raggedright\arraybackslash}X}
        \toprule
        \textbf{Rolle} & \textbf{User Story} \\
        \midrule
        Leser:in & Als Leser:in möchte ich durch kurze Zusammenfassungen wischen, um in kurzer Zeit einen Überblick über die aktuellen Nachrichten zu bekommen. \\
        Leser:in & Als Leser:in möchte ich einen Artikel mit einem Tippen vollständig öffnen, um mehr über ein Thema zu erfahren. \\
        Leser:in & Als Leser:in möchte ich den Feed nach Kategorien filtern, um nur Themen zu sehen, die mich gerade interessieren. \\
        Leser:in & Als Leser:in möchte ich, dass der Feed mit der Zeit mehr Artikel zeigt, die zu meinen Interessen passen, ohne dass ich mich anmelden muss. \\
        Leser:in & Als Leser:in möchte ich Artikel liken, kommentieren und teilen können. \\
        Redakteur:in & Als Redakteur:in möchte ich einen Artikel anlegen und automatisch Teaser, Kernpunkte und Tags erhalten, um keine Kurzfassung selbst schreiben zu müssen. \\
        Redakteur:in & Als Redakteur:in möchte ich generierte Texte korrigieren und frühere Fassungen wiederherstellen können. \\
        Redakteur:in & Als Redakteur:in möchte ich Artikel aus einem bestehenden \gls{rss}-Feed übernehmen, um sie nicht abtippen zu müssen. \\
        Redakteur:in & Als Redakteur:in möchte ich sehen, wie oft ein Artikel angesehen und wie lange er gelesen wurde. \\
        Entwickler:in & Als Entwickler:in möchte ich den Feed mit einer einzigen Komponente in mein Nuxt-Projekt einbauen. \\
        Entwickler:in & Als Entwickler:in möchte ich Farben und Schriften des Feeds an meine Website anpassen. \\
        Entwickler:in & Als Entwickler:in möchte ich selbst entscheiden, was beim Öffnen eines Artikels passiert. \\
        \bottomrule
    \end{tabularx}
    \caption{Auswahl der User Stories je Rolle}
    \label{tab:req:stories}
\end{table}

Aus den User Stories ergeben sich zwei zentrale Anwendungsfälle, die das ganze System durchlaufen.
Im ersten legt ein:e Redakteur:in einen Artikel an.
Das System übernimmt Titel, Text und Kategorie, erzeugt im Hintergrund Teaser, Kernpunkte und Tags und zeigt den Fortschritt im Admin-\gls{cms} an.
Nach der Prüfung wird der Artikel veröffentlicht und erscheint beim nächsten Abruf im Feed.
Ändert sich später der Artikeltext, wird der Teaser neu erzeugt, während eine von Hand geänderte Fassung erhalten bleibt.

Im zweiten Anwendungsfall öffnet eine Person die Website eines Mediums auf dem Smartphone.
Der Feed lädt die ersten Artikel, zeigt die erste Karte und lädt weitere Karten nach, während gewischt wird.
Aus dem Verhalten, etwa wie lange eine Karte betrachtet oder ob ein Artikel geöffnet wird, lernt der Feed Interessen und ordnet die folgenden Karten neu.
Tippt die Person auf \glqq Weiterlesen\grqq{}, meldet der Feed die Kennung des Artikels an die Website, die den vollständigen Text anzeigt.

Beide Anwendungsfälle zeigen, dass Generierung und Darstellung nur über die Schnittstelle miteinander verbunden sind.
Die genauen Anforderungen an beide Seiten folgen im nächsten Abschnitt.

\section{Funktionale Anforderungen}\label{sec:req:functional}
\setauthor{\firstauthor}

Tabelle~\ref{tab:req:functional} listet die funktionalen Pflichtanforderungen.
Die Kennungen mit F, G und E stammen aus dem Briefing, jene mit D aus dem Diplomarbeitsantrag.
Anforderungen, die in beiden Dokumenten vorkommen, sind nur einmal angeführt.

\begin{table}
    \centering
    \small
    \begin{tabularx}{\textwidth}{l >{\raggedright\arraybackslash}X}
        \toprule
        \textbf{ID} & \textbf{Anforderung} \\
        \midrule
        F1 & Vertikal scrollbarer Short-Form-Feed mit einer Karte pro Artikel \\
        F2 & Anzeige von KI-generierten Kurzinhalten (Teaser) auf den Karten \\
        F3 & Übergabe einer eindeutigen Kennung an die Host-Website, wenn ein Artikel geöffnet wird \\
        G1 & Generierung der Kurzinhalte aus strukturierten Artikeldaten \\
        G2 & Dauerhafte Speicherung der generierten Inhalte \\
        E1 & Einbindung des Feeds als Nuxt-Komponente in bestehende Nuxt-Projekte \\
        E2 & Erkennung neuer und geänderter Inhalte im Feed \\
        E3 & Bereitstellung der Inhalte über eine \gls{rest}-Schnittstelle \\
        D1 & Reproduzierbare und konfigurierbare KI-Generierung \\
        D2 & Asynchrone Verarbeitung neuer und geänderter Inhalte \\
        D3 & Versionierung der generierten Inhalte \\
        D4 & Unabhängigkeit von einem bestimmten \gls{cms} \\
        \bottomrule
    \end{tabularx}
    \caption{Funktionale Pflichtanforderungen}
    \label{tab:req:functional}
\end{table}

Einige dieser Anforderungen sind im Briefing bewusst knapp formuliert und mussten für die Umsetzung präzisiert werden.

\textbf{F3} legt fest, dass der Feed beim Öffnen eines Artikels nur eine Kennung übergibt.
Der Feed zeigt also nicht selbst eine eigene Artikelseite an, sondern meldet über ein Ereignis die Artikel-ID an die einbindende Website.
Damit bleibt die Website für Navigation, Layout und gegebenenfalls Bezahlschranken zuständig.

\textbf{G1} wurde so präzisiert, dass aus Titel, Text und Kategorie eines Artikels ein Teaser mit höchstens drei Sätzen, drei Kernpunkte und ein bis zwei Tags entstehen.
Die Tags stammen aus einer festen Liste von Unterkategorien je Hauptkategorie.
Freie Schlagwörter wie Personen- oder Firmennamen würden zu sehr vielen Tags mit jeweils nur einem Artikel führen und wären für die Filterung im Feed unbrauchbar.
Als Hauptkategorien sind Politik, Wirtschaft, Sport, Technologie, Kultur und Chronik festgelegt.
Zusätzlich zum Text entsteht auf Wunsch ein Titelbild im Hochformat.

\textbf{E2} verlangt, dass der Feed neue oder geänderte Inhalte erkennt, ohne dass die Seite neu geladen werden muss.
Das Briefing lässt offen, ob das über regelmäßige Abfragen oder über eine Push-Verbindung geschieht.

\textbf{D1} verlangt, dass derselbe Artikel mit denselben Einstellungen dasselbe Ergebnis liefert.
Bei Sprachmodellen ist das nicht selbstverständlich, weil die Ausgabe standardmäßig zufällig gezogen wird.
Modell, Temperatur und Seed müssen deshalb einstellbar sein und bei jedem Ergebnis mitgespeichert werden.

\textbf{D3} geht über das bloße Speichern hinaus.
Jede generierte oder von Hand geänderte Fassung soll als eigene, unveränderliche Version erhalten bleiben und wiederhergestellt werden können.
Das ist die Grundlage dafür, dass Redakteur:innen Ergebnisse vergleichen und die in Abschnitt~\ref{sec:related:summarization} beschriebenen Fehler eines Sprachmodells nachvollziehen können.

\textbf{D4} bedeutet, dass Artikel aus beliebigen Quellen ins System gelangen können.
Neben der direkten Anlage über die Schnittstelle ist dafür ein Import aus \gls{rss}- und Atom-Feeds vorgesehen, sodass bestehende Nachrichtenquellen ohne eigene Anpassung angebunden werden können.

Neben den Pflichtanforderungen nennt das Briefing mehrere Erweiterungsideen, die in Tabelle~\ref{tab:req:optional} zusammengefasst sind.
Sie fließen nicht in die Bewertung der Pflichtanforderungen ein, wurden aber bei der Architektur von Anfang an berücksichtigt.

\begin{table}
    \centering
    \small
    \begin{tabularx}{\textwidth}{l >{\raggedright\arraybackslash}X}
        \toprule
        \textbf{ID} & \textbf{Optionale Erweiterung} \\
        \midrule
        O1 & Personalisierung der Reihenfolge im Feed \\
        O2 & Personalisierung der Inhalte, etwa Teaser-Varianten je Zielgruppe \\
        O3 & Admin-Dashboard auf Basis von Nuxt~4 \\
        O4 & Weitergabe von Nutzungsdaten über Client-Ereignisse oder Webhooks \\
        O5 & Audiofassung der Inhalte \\
        \bottomrule
    \end{tabularx}
    \caption{Optionale Erweiterungen laut Briefing}
    \label{tab:req:optional}
\end{table}

Die funktionalen Anforderungen beschreiben, was das System leisten soll.
Wie schnell, wie sicher und wie gut wartbar es dabei sein muss, legen die nicht-funktionalen Anforderungen im folgenden Abschnitt fest.

\section{Nicht-funktionale Anforderungen}
\setauthor{\secondauthor}

Neben den technischen Rahmenbedingungen aus Tabelle~\ref{tab:req:tech} enthält der Diplomarbeitsantrag mehrere Ziele, die die Qualität des Systems betreffen.
Sie sind in Tabelle~\ref{tab:req:nonfunctional} zusammengefasst und werden anschließend erläutert.

\begin{table}
    \centering
    \small
    \begin{tabularx}{\textwidth}{l >{\raggedright\arraybackslash}X}
        \toprule
        \textbf{ID} & \textbf{Anforderung} \\
        \midrule
        S1 & Responsive Oberfläche, ausgelegt für Smartphones im Hochformat \\
        S2 & Performantes Rendering und schrittweises Nachladen der Inhalte \\
        S3 & Konfigurierbare Darstellung und Gestaltung des Feeds \\
        S4 & Ereignisse für Tracking und Analytics \\
        S5 & Technische Dokumentation zur Einbindung \\
        D5 & Performance des Backends durch Caching und Warteschlangen \\
        D6 & Dokumentation von Schnittstellen und Architektur \\
        \bottomrule
    \end{tabularx}
    \caption{Nicht-funktionale Anforderungen}
    \label{tab:req:nonfunctional}
\end{table}

\textbf{Benutzbarkeit auf dem Smartphone (S1).}
Der Feed muss auf Bildschirmen ab einer Breite von 375 Pixeln vollständig bedienbar sein.
Jede Karte füllt die Höhe des Bildschirms aus und rastet beim Wischen ein, sodass immer genau ein Artikel sichtbar ist.
Bedienelemente wie Like oder Teilen müssen mit dem Daumen erreichbar sein.

\textbf{Ladezeit und Speicherbedarf (S2, D5).}
Ein Feed kann sehr viele Artikel enthalten, und jede Karte trägt ein großes Bild.
Würden beim Start alle Artikel geladen und alle Karten gezeichnet, stiegen Datenmenge und Speicherbedarf mit jedem neuen Artikel.
Der Feed soll daher Artikel in Seiten laden und nur die Karten rendern, die sich in der Nähe der sichtbaren Karte befinden.
Auf der Seite des Backends darf eine langsame Generierung keine Anfrage an den Feed blockieren.
Generierung und Auslieferung müssen deshalb getrennt sein, und häufig abgefragte Daten sollen zwischengespeichert werden.
Da der Feed regelmäßig auf Änderungen prüft (E2), soll eine unveränderte Antwort ohne erneute Übertragung des Inhalts beantwortet werden können.

\textbf{Anpassbarkeit (S3).}
Websites, die den Feed einbinden, haben eigene Farben und Schriften.
Das Aussehen des Feeds muss daher ohne Änderung am Code des Packages anpassbar sein.
Einzelne Bereiche einer Karte sollen sich zusätzlich durch eigene Inhalte ersetzen lassen.

\textbf{Messbarkeit (S4).}
Für Redaktionen ist es wichtig zu wissen, welche Artikel im Feed funktionieren.
Der Feed muss deshalb Ereignisse wie das erste Sichtbarwerden einer Karte, das Öffnen eines Artikels oder ein Like nach außen melden.
Die einbindende Website kann diese Ereignisse an ein eigenes Analysewerkzeug weitergeben.

\textbf{Dokumentation (S5, D6).}
Die Schnittstelle soll vollständig und maschinenlesbar beschrieben sein, sodass andere Entwickler:innen sie ohne Blick in den Quellcode nutzen können.
Für das Package ist eine Anleitung nötig, die alle Schritte der Einbindung, die Konfiguration und alle Ereignisse beschreibt.

\textbf{Sicherheit und Datenschutz.}
Diese Punkte sind im Briefing nicht ausdrücklich genannt, ergeben sich aber aus dem Einsatz auf öffentlichen Websites.
Schreibende Zugriffe auf die Schnittstelle, etwa das Anlegen oder Löschen von Artikeln, dürfen nur mit einem Schlüssel möglich sein.
Der Feed selbst muss ohne Anmeldung funktionieren.
Das Interessenprofil für die Personalisierung soll das Gerät der Leser:innen nicht verlassen, an den Server gehen nur anonyme Ereignisse für die Statistik.
Eine weitere Vorgabe des Projektteams betrifft die Kosten: Text- und Bildgenerierung sollen lokal auf eigener Hardware laufen können, damit keine laufenden Kosten pro Artikel bei einem externen Anbieter entstehen.

Mit diesen Anforderungen steht fest, was das System leisten muss.
Wie die einzelnen Bausteine aufgebaut sind, um sie zu erfüllen, zeigt das folgende Kapitel.
